\setcounter{chapter}{8}
\pagestyle{fancy}

\chapter{矩阵综合}

本章主要作为一个归纳章节，大部分内容在前文已经以定理、命题或例题的形式出现。
这里按照原讲义的组织重新汇总，便于复习矩阵分解与常见特殊矩阵。


\section{矩阵分解}

\subsection{满秩分解}

\begin{theorem}[\markstar]
设
\[
A\in M_{m\times n}(F),
\qquad
r(A)=r.
\]
则存在
\[
F\in M_{m\times r}(F),
\qquad
G\in M_{r\times n}(F),
\]
使
\[
A=FG,
\qquad
r(F)=r(G)=r.
\]
这样的分解称为 $A$ 的满秩分解。
\end{theorem}

具体求解时，可以先对 $A$ 作初等行变换得到
\[
CA=
\begin{pmatrix}
G\\
O
\end{pmatrix},
\]
再从 $C^{-1}$ 中取对应的前 $r$ 列组成 $F$；也可以直接选取 $A$ 的
$r$ 个线性无关列组成 $F$，然后把其余各列写成这些列的线性组合，
相应坐标组成 $G$。


\begin{example}
求下列矩阵的满秩分解：
\begin{enumerate}
    \item
    \[
    A=
    \begin{pmatrix}
    1&2&3&0\\
    0&2&1&-1\\
    1&0&2&1
    \end{pmatrix};
    \]
    \item
    \[
    A=
    \begin{pmatrix}
    1&3&2&1&4\\
    2&6&1&0&7\\
    3&9&3&1&11
    \end{pmatrix}.
    \]
\end{enumerate}
\end{example}
\exampleblank

\begin{solution}
\begin{enumerate}
    \item 对 $A$ 行化简得
    \[
    A\sim
    \begin{pmatrix}
    1&0&2&1\\
    0&1&\dfrac12&-\dfrac12\\
    0&0&0&0
    \end{pmatrix},
    \]
    因而 $r(A)=2$，第一、二列可作为列空间的一组基。取
    \[
    F=
    \begin{pmatrix}
    1&2\\
    0&2\\
    1&0
    \end{pmatrix},
    \qquad
    G=
    \begin{pmatrix}
    1&0&2&1\\
    0&1&\dfrac12&-\dfrac12
    \end{pmatrix}.
    \]
    则
    \[
    \boxed{A=FG}.
    \]

    \item 行化简得
    \[
    A\sim
    \begin{pmatrix}
    1&3&0&-\dfrac13&\dfrac{10}{3}\\
    0&0&1&\dfrac23&\dfrac13\\
    0&0&0&0&0
    \end{pmatrix},
    \]
    所以 $r(A)=2$，第一、三列线性无关。取
    \[
    F=
    \begin{pmatrix}
    1&2\\
    2&1\\
    3&3
    \end{pmatrix},
    \qquad
    G=
    \begin{pmatrix}
    1&3&0&-\dfrac13&\dfrac{10}{3}\\
    0&0&1&\dfrac23&\dfrac13
    \end{pmatrix}.
    \]
    则
    \[
    \boxed{A=FG}.
    \]
\end{enumerate}
\end{solution}


\subsection{平方根分解}

\begin{theorem}[\marktwostar]
设 $A$ 为实对称半正定矩阵，则存在唯一实对称半正定矩阵
\[
X=A^{1/2}
\]
满足
\[
X^2=A.
\]
\end{theorem}


\begin{example}
求下列矩阵的半正定平方根：
\begin{enumerate}
    \item
    \[
    A_1=
    \begin{pmatrix}
    4&2\\
    2&1
    \end{pmatrix};
    \]
    \item
    \[
    A_2=
    \begin{pmatrix}
    9&3&3\\
    3&1&1\\
    3&1&1
    \end{pmatrix}.
    \]
\end{enumerate}
\end{example}
\exampleblank

\begin{solution}
\begin{enumerate}
    \item
    \[
    A_1=
    \begin{pmatrix}2\\1\end{pmatrix}
    \begin{pmatrix}2&1\end{pmatrix},
    \qquad
    A_1^2=5A_1.
    \]
    因此
    \[
    \boxed{
    A_1^{1/2}
    =
    \dfrac1{\sqrt5}
    \begin{pmatrix}
    4&2\\
    2&1
    \end{pmatrix}.
    }
    \]

    \item
    \[
    A_2=
    \begin{pmatrix}3\\1\\1\end{pmatrix}
    \begin{pmatrix}3&1&1\end{pmatrix},
    \qquad
    A_2^2=11A_2.
    \]
    因此
    \[
    \boxed{
    A_2^{1/2}
    =
    \dfrac1{\sqrt{11}}
    \begin{pmatrix}
    9&3&3\\
    3&1&1\\
    3&1&1
    \end{pmatrix}.
    }
    \]
\end{enumerate}
\end{solution}


\subsection{Schur 分解}

\begin{theorem}[\markstar]
设 $A\in M_n(\mathbb C)$。则存在酉矩阵 $U$，使
\[
U^HAU=T,
\]
其中 $T$ 为上三角矩阵，其对角元恰为 $A$ 的全部特征值。
这称为复 Schur 分解。
\end{theorem}

\begin{theorem}[\markstar]
设 $A\in M_n(\mathbb R)$。则存在正交矩阵 $Q$，使
\[
Q^{\mathsf T}AQ
\]
为实准上三角矩阵，其对角块是一阶实块或由一对共轭复特征值对应的
二阶实块。这称为实 Schur 分解。
\end{theorem}


\begin{example}
求下列矩阵的实 Schur 分解：
\begin{enumerate}
    \item
    \[
    A_1=
    \begin{pmatrix}
    1&2\\
    2&1
    \end{pmatrix};
    \]
    \item
    \[
    A_2=
    \begin{pmatrix}
    0&1&0\\
    -1&0&0\\
    0&0&1
    \end{pmatrix}.
    \]
\end{enumerate}
\end{example}
\exampleblank

\begin{solution}
\begin{enumerate}
    \item $A_1$ 为实对称矩阵。取
    \[
    Q=
    \dfrac1{\sqrt2}
    \begin{pmatrix}
    1&1\\
    1&-1
    \end{pmatrix},
    \]
    则
    \[
    \boxed{
    Q^{\mathsf T}A_1Q
    =
    \begin{pmatrix}
    3&0\\
    0&-1
    \end{pmatrix}.
    }
    \]

    \item $A_2$ 已经是实 Schur 标准形式：
    \[
    A_2=
    \begin{pmatrix}
    \boxed{\begin{matrix}0&1\\-1&0\end{matrix}}&0\\
    0&1
    \end{pmatrix}.
    \]
    因此可直接取
    \[
    \boxed{Q=I_3}.
    \]
\end{enumerate}
\end{solution}


\begin{example}
求下列矩阵的复 Schur 分解：
\begin{enumerate}
    \item
    \[
    A_1=
    \begin{pmatrix}
    1&i\\
    -i&1
    \end{pmatrix};
    \]
    \item
    \[
    A_2=
    \begin{pmatrix}
    i&1\\
    0&-i
    \end{pmatrix}.
    \]
\end{enumerate}
\end{example}
\exampleblank

\begin{solution}
\begin{enumerate}
    \item $A_1$ 为 Hermite 矩阵，特征值为 $0,2$。取单位特征向量
    \[
    u_1=\dfrac1{\sqrt2}
    \begin{pmatrix}
    -i\\1
    \end{pmatrix},
    \qquad
    u_2=\dfrac1{\sqrt2}
    \begin{pmatrix}
    i\\1
    \end{pmatrix},
    \]
    并令
    \[
    U=(u_1,u_2).
    \]
    则
    \[
    \boxed{
    U^HA_1U=\diag(0,2).
    }
    \]

    \item $A_2$ 本身已经是上三角矩阵，所以可取
    \[
    \boxed{U=I_2,\qquad T=A_2}.
    \]
\end{enumerate}
\end{solution}


\subsection{谱分解}

\begin{theorem}[\marktwostar]
设 $A\in M_n(\mathbb C)$ 为正规矩阵，即
\[
A^HA=AA^H.
\]
则存在酉矩阵 $U$ 使
\[
U^HAU
=
\diag(\lambda_1,\ldots,\lambda_n).
\]
\end{theorem}


\begin{example}
求下列正规矩阵的谱分解：
\begin{enumerate}
    \item
    \[
    A_1=
    \begin{pmatrix}
    1&1\\
    -1&1
    \end{pmatrix};
    \]
    \item
    \[
    A_2=
    \begin{pmatrix}
    0&1&0\\
    1&0&0\\
    0&0&1
    \end{pmatrix}.
    \]
\end{enumerate}
\end{example}
\exampleblank

\begin{solution}
\begin{enumerate}
    \item 特征值为
    \[
    1+i,\qquad 1-i.
    \]
    取
    \[
    u_1=
    \dfrac1{\sqrt2}
    \begin{pmatrix}1\\i\end{pmatrix},
    \qquad
    u_2=
    \dfrac1{\sqrt2}
    \begin{pmatrix}1\\-i\end{pmatrix},
    \]
    令
    \[
    U=(u_1,u_2),
    \]
    则
    \[
    \boxed{
    U^HA_1U
    =
    \diag(1+i,1-i).
    }
    \]

    \item 特征值为 $1,-1,1$。取
    \[
    u_1=
    \dfrac1{\sqrt2}
    \begin{pmatrix}1\\1\\0\end{pmatrix},
    \quad
    u_2=
    \dfrac1{\sqrt2}
    \begin{pmatrix}1\\-1\\0\end{pmatrix},
    \quad
    u_3=
    \begin{pmatrix}0\\0\\1\end{pmatrix}.
    \]
    则
    \[
    U=(u_1,u_2,u_3)
    \]
    为正交矩阵，并且
    \[
    \boxed{
    U^{\mathsf T}A_2U=\diag(1,-1,1).
    }
    \]
\end{enumerate}
\end{solution}


\subsection{奇异值分解}

\begin{theorem}[\markstar]
设
\[
A\in M_{m\times n}(\mathbb C),
\qquad
r(A)=r.
\]
则存在酉矩阵 $U\in M_m(\mathbb C)$、$V\in M_n(\mathbb C)$，使
\[
U^HAV
=
\begin{pmatrix}
\Sigma_r&O\\
O&O
\end{pmatrix},
\]
其中
\[
\Sigma_r=\diag(\sigma_1,\ldots,\sigma_r),
\qquad
\sigma_1\ge\cdots\ge\sigma_r>0,
\]
而 $\sigma_i^2$ 是 $A^HA$ 的非零特征值。
\end{theorem}


\begin{example}
求下列矩阵的奇异值分解：
\begin{enumerate}
    \item
    \[
    A_1=
    \begin{pmatrix}
    1&2&0\\
    3&6&0
    \end{pmatrix};
    \]
    \item
    \[
    A_2=
    \begin{pmatrix}
    1&0&1\\
    0&1&1\\
    0&0&0
    \end{pmatrix}.
    \]
\end{enumerate}
\end{example}
\exampleblank

\begin{solution}
\begin{enumerate}
    \item $A_1$ 的秩为 $1$，唯一非零奇异值为
    \[
    \sigma_1=5\sqrt2.
    \]
    取
    \[
    U=
    \dfrac1{\sqrt{10}}
    \begin{pmatrix}
    1&-3\\
    3&1
    \end{pmatrix},
    \]
    \[
    V=
    \begin{pmatrix}
    \dfrac1{\sqrt5}&-\dfrac2{\sqrt5}&0\\[4pt]
    \dfrac2{\sqrt5}&\dfrac1{\sqrt5}&0\\
    0&0&1
    \end{pmatrix}.
    \]
    则
    \[
    \boxed{
    U^{\mathsf T}A_1V
    =
    \begin{pmatrix}
    5\sqrt2&0&0\\
    0&0&0
    \end{pmatrix}.
    }
    \]

    \item $A_2^{\mathsf T}A_2$ 的特征值为
    \[
    3,\quad1,\quad0,
    \]
    所以奇异值为
    \[
    \sqrt3,\quad1.
    \]
    可取
    \[
    U=
    \begin{pmatrix}
    \dfrac1{\sqrt2}&\dfrac1{\sqrt2}&0\\[4pt]
    \dfrac1{\sqrt2}&-\dfrac1{\sqrt2}&0\\
    0&0&1
    \end{pmatrix},
    \]
    \[
    V=
    \begin{pmatrix}
    \dfrac1{\sqrt6}&\dfrac1{\sqrt2}&-\dfrac1{\sqrt3}\\[4pt]
    \dfrac1{\sqrt6}&-\dfrac1{\sqrt2}&-\dfrac1{\sqrt3}\\[4pt]
    \dfrac2{\sqrt6}&0&\dfrac1{\sqrt3}
    \end{pmatrix}.
    \]
    则
    \[
    \boxed{
    U^{\mathsf T}A_2V
    =
    \diag(\sqrt3,1,0).
    }
    \]
\end{enumerate}
\end{solution}


\subsection{极分解}

\begin{theorem}[\markstar]
设 $A\in M_n(\mathbb C)$。则存在酉矩阵 $U$ 与 Hermite 半正定矩阵
$H_1,H_2$，使
\[
A=H_1U=UH_2,
\]
其中
\[
H_1=(AA^H)^{1/2},
\qquad
H_2=(A^HA)^{1/2}.
\]
半正定因子 $H_1,H_2$ 唯一；当 $A$ 可逆时，酉因子也唯一。
\end{theorem}


\begin{example}
求下列矩阵的极分解：
\begin{enumerate}
    \item
    \[
    A_1=
    \begin{pmatrix}
    1&1\\
    0&1
    \end{pmatrix};
    \]
    \item
    \[
    A_2=
    \begin{pmatrix}
    1&1&0\\
    0&1&1\\
    1&0&1
    \end{pmatrix}.
    \]
\end{enumerate}
\end{example}
\exampleblank

\begin{solution}
\begin{enumerate}
    \item 有
    \[
    A_1A_1^{\mathsf T}
    =
    \begin{pmatrix}
    2&1\\
    1&1
    \end{pmatrix},
    \qquad
    A_1^{\mathsf T}A_1
    =
    \begin{pmatrix}
    1&1\\
    1&2
    \end{pmatrix}.
    \]
    它们的正定平方根分别为
    \[
    H_1=
    \dfrac1{\sqrt5}
    \begin{pmatrix}
    3&1\\
    1&2
    \end{pmatrix},
    \qquad
    H_2=
    \dfrac1{\sqrt5}
    \begin{pmatrix}
    2&1\\
    1&3
    \end{pmatrix}.
    \]
    取
    \[
    U=
    \dfrac1{\sqrt5}
    \begin{pmatrix}
    2&1\\
    -1&2
    \end{pmatrix},
    \]
    则
    \[
    \boxed{A_1=H_1U=UH_2}.
    \]

    \item 计算
    \[
    A_2A_2^{\mathsf T}
    =
    A_2^{\mathsf T}A_2
    =
    I_3+J_3.
    \]
    该矩阵在 $(1,1,1)^{\mathsf T}$ 方向上的特征值为 $4$，
    在其正交补上的特征值为 $1$，因此
    \[
    H_1=H_2=
    I_3+\dfrac13J_3
    =
    \dfrac13
    \begin{pmatrix}
    4&1&1\\
    1&4&1\\
    1&1&4
    \end{pmatrix}.
    \]
    相应正交因子为
    \[
    U=
    \dfrac13
    \begin{pmatrix}
    2&2&-1\\
    -1&2&2\\
    2&-1&2
    \end{pmatrix}.
    \]
    故
    \[
    \boxed{A_2=H_1U=UH_2}.
    \]
\end{enumerate}
\end{solution}


\subsection{QR 分解}

\begin{theorem}[\markstar]
设 $A$ 为实可逆矩阵，则存在唯一正交矩阵 $Q$ 与主对角元均为正数的
上三角矩阵 $R$，使
\[
A=QR.
\]
这正是 Gram--Schmidt 正交化的矩阵形式。
\end{theorem}


\begin{example}
求下列矩阵的 QR 分解：
\begin{enumerate}
    \item
    \[
    A_1=
    \begin{pmatrix}
    1&1\\
    1&2
    \end{pmatrix};
    \]
    \item
    \[
    A_2=
    \begin{pmatrix}
    1&1&1\\
    1&2&3\\
    1&3&6
    \end{pmatrix}.
    \]
\end{enumerate}
\end{example}
\exampleblank

\begin{solution}
\begin{enumerate}
    \item 对列向量作 Gram--Schmidt 正交化，得到
    \[
    Q=
    \dfrac1{\sqrt2}
    \begin{pmatrix}
    1&-1\\
    1&1
    \end{pmatrix},
    \]
    \[
    R=
    \begin{pmatrix}
    \sqrt2&\dfrac3{\sqrt2}\\[4pt]
    0&\dfrac1{\sqrt2}
    \end{pmatrix}.
    \]
    因此
    \[
    \boxed{A_1=QR}.
    \]

    \item 依次取
    \[
    q_1=\dfrac1{\sqrt3}(1,1,1)^{\mathsf T},
    \]
    \[
    q_2=\dfrac1{\sqrt2}(-1,0,1)^{\mathsf T},
    \]
    \[
    q_3=\dfrac1{\sqrt6}(1,-2,1)^{\mathsf T}.
    \]
    令
    \[
    Q=(q_1,q_2,q_3),
    \]
    则
    \[
    R=
    \begin{pmatrix}
    \sqrt3&2\sqrt3&\dfrac{10}{\sqrt3}\\[4pt]
    0&\sqrt2&\dfrac5{\sqrt2}\\[4pt]
    0&0&\dfrac1{\sqrt6}
    \end{pmatrix}.
    \]
    因而
    \[
    \boxed{A_2=QR}.
    \]
\end{enumerate}
\end{solution}


\newpage

\section{特殊矩阵}

本节按照原讲义要求，对前文已经多次出现的特殊矩阵作集中归纳。

\begin{example}[\marktwostar]
分析下列矩阵的主要性质，包括特征值、标准形以及常见等价刻画：
\begin{enumerate}
    \item 实对称（Hermite）矩阵；
，    \item 实反称（反 Hermite）矩阵；
    \item 正交（酉）矩阵；
    \item 正规矩阵；
    \item 幂零矩阵；
    \item 幂等矩阵；
    \item 秩 $1$ 矩阵；
    \item 镜像矩阵；
    \item 对合矩阵；
    \item 循环矩阵。
\end{enumerate}
\end{example}
\exampleblank

\begin{solution}
\begin{enumerate}
    \item \textbf{实对称 / Hermite 矩阵。}
    实对称矩阵满足
    \[
    A^{\mathsf T}=A,
    \]
    Hermite 矩阵满足
    \[
    A^H=A.
    \]
    特征值全为实数；属于不同特征值的特征向量正交；存在正交或酉矩阵
    $Q$ 使
    \[
    Q^HAQ=\diag(\lambda_1,\ldots,\lambda_n).
    \]
    因而一定可对角化。

    \item \textbf{实反称 / 反 Hermite 矩阵。}
    分别满足
    \[
    A^{\mathsf T}=-A,
    \qquad
    A^H=-A.
    \]
    反 Hermite 矩阵的特征值均为纯虚数或 $0$。
    实反称矩阵的非零特征值成对为
    \[
    \pm i\mu_j.
    \]
    实正交相似标准形由二维块
    \[
    \begin{pmatrix}
    0&\mu_j\\
    -\mu_j&0
    \end{pmatrix}
    \]
    与零块组成。

    \item \textbf{正交 / 酉矩阵。}
    分别满足
    \[
    A^{\mathsf T}A=I,
    \qquad
    A^HA=I.
    \]
    等价地，保持内积与长度。全部复特征值满足
    \[
    |\lambda|=1.
    \]
    正交、酉矩阵都是正规矩阵，因此在复数域上可酉对角化。
    实正交矩阵在实数域上的标准形由 $\pm1$ 与二维旋转块组成。

    \item \textbf{正规矩阵。}
    满足
    \[
    A^HA=AA^H.
    \]
    其核心等价刻画是：存在酉矩阵 $U$ 使
    \[
    U^HAU
    \]
    为对角矩阵。Hermite、反 Hermite、酉矩阵都是正规矩阵。

    \item \textbf{幂零矩阵。}
    存在正整数 $k$ 使
    \[
    A^k=O.
    \]
    全部特征值均为 $0$；特征多项式为 $\lambda^n$；
    Jordan 标准形只含零特征值的 Jordan 块。
    若最小的 $k$ 满足 $A^k=0$，则最大 Jordan 块的阶数为 $k$。

    \item \textbf{幂等矩阵。}
    满足
    \[
    A^2=A.
    \]
    最小多项式整除
    \[
    x(x-1),
    \]
    因而一定可对角化，特征值只能为 $0,1$。并且
    \[
    \boxed{r(A)=\operatorname{tr}A}.
    \]
    几何上它对应某个直和分解下的投影。

    \item \textbf{秩 $1$ 矩阵。}
    非零秩一矩阵可写成
    \[
    A=uv^{\mathsf T}.
    \]
    有
    \[
    A^2=(v^{\mathsf T}u)A=(\operatorname{tr}A)A.
    \]
    因而其特征值中至多有一个非零值，且该非零值为
    \[
    \operatorname{tr}A.
    \]
    若 $\operatorname{tr}A\neq0$，则 $A$ 可对角化；
    若 $\operatorname{tr}A=0$，则 $A^2=0$。

    \item \textbf{镜像矩阵。}
    对单位向量 $w$，
    \[
    H=I-2ww^{\mathsf T}.
    \]
    它满足
    \[
    H^{\mathsf T}=H,\qquad H^2=I,
    \]
    因而既对称又正交。其特征值为
    \[
    -1\quad(1\text{ 重}),\qquad
    1\quad(n-1\text{ 重}),
    \]
    且
    \[
    \det H=-1.
    \]

    \item \textbf{对合矩阵。}
    满足
    \[
    A^2=I.
    \]
    在特征不为 $2$ 的数域上，最小多项式整除
    \[
    (x-1)(x+1),
    \]
    无重根，因此 $A$ 可对角化，特征值只可能为 $\pm1$。
    若再满足 $A^{\mathsf T}=A$，则它是正交矩阵。

    \item \textbf{循环矩阵。}
    若
    \[
    A=
    \operatorname{circ}(a_0,a_1,\ldots,a_{n-1}),
    \]
    令
    \[
    f(x)=a_0+a_1x+\cdots+a_{n-1}x^{n-1},
    \]
    则其特征值为
    \[
    f(\omega_k),
    \]
    其中 $\omega_k$ 遍历全部 $n$ 次单位根。
    离散 Fourier 矩阵可同时对角化所有循环矩阵，因此循环矩阵彼此可交换，
    且
    \[
    \det A=\prod\limits_{k=1}^{n}f(\omega_k).
    \]
\end{enumerate}
\end{solution}
